\chapter{Percutaneous Tumor Ablation}

\section*{Overview}
Percutaneous tumor ablation is a minimally invasive technique in which image guidance is used to place a needle or probe directly into a tumor to destroy it by extreme cold (cryoablation) or heat (radiofrequency or microwave ablation). It is performed through the skin without open surgery, typically as an outpatient or short-stay procedure.

\section*{Alternative Names}
Cryoablation, cryotherapy, radiofrequency ablation (RFA), microwave ablation (MWA), thermal ablation, tumor ablation, image-guided tumor ablation

\section*{Principle}
All three modalities rely on a probe positioned within the tumor under computed tomography, ultrasound, or magnetic resonance imaging guidance. Cryoablation destroys cells by rapid freeze--thaw cycles that cause ice-crystal formation, osmotic injury, and microvascular thrombosis. Radiofrequency and microwave ablation generate frictional and molecular heat, respectively, raising tissue temperature above 50--60\textdegree C to produce coagulative necrosis. The goal is complete necrosis of the target lesion while sparing adjacent structures.

\section*{History}
Percutaneous ablation was developed in the 1990s and expanded rapidly as image-guidance technology improved. Cryoablation predates heat-based methods, with modern imaging-compatible systems now allowing real-time monitoring of the ice ball.

\section*{Why It Is Done}
Percutaneous ablation is used to treat primary and metastatic tumors of the liver, kidney, lung, bone, and adrenal glands when surgery is not feasible because of patient comorbidities, tumor number or location, or poor surgical candidacy. It is also used for pain palliation of osseous metastases and to control tumor burden in patients awaiting transplant or other systemic therapy. In selected cases it offers local control comparable to surgery with less morbidity and shorter recovery.

\section*{Is This a Primary or Secondary Test or Procedure}
It is a primary therapeutic procedure for appropriately selected small tumors. It is a secondary or alternative option when resection, transplant, or stereotactic radiation is contraindicated or declined.

\section*{How It Is Performed}
The patient is positioned, and the target tumor is localized with CT, ultrasound, or MRI. After sterile preparation and local or general anesthesia, one or more ablation probes are inserted percutaneously into the lesion. Energy is delivered in cycles to create a margin of necrosis around the tumor, and for cryoablation the ice ball is monitored by real-time imaging. Multiple ablations may be performed during the same session to cover larger or irregular tumors.

\section*{Preparation}
Preprocedural imaging and laboratory evaluation, including coagulation studies, are required, and anticoagulant or antiplatelet therapy is typically held before the procedure. Patients fast for several hours and are instructed about anesthesia and postprocedural observation. Written informed consent is obtained after a discussion of risks, benefits, and alternatives.

\section*{Contraindications and Precautions}
Uncorrectable coagulopathy, active infection, and inability to safely reach the tumor are relative contraindications. Lesions adjacent to major vessels, bile ducts, bowel, or other critical structures carry an increased risk of collateral injury. Tumor size, location, and number must be assessed because larger or exophytic tumors have higher incomplete-ablation rates.

\section*{Risks}
Risks include bleeding, infection, pneumothorax or pleural effusion for lung lesions, bile duct or bowel injury, thermal injury to adjacent organs, and tumor seeding along the needle track. Pain and a postablation syndrome of fever, malaise, and nausea may occur transiently.

\section*{Aftercare and Recovery}
Patients are monitored in a recovery area for several hours, and imaging is performed to assess for complications such as bleeding or pneumothorax. Most patients return home the same day or after a short overnight stay and resume normal activity within a few days. Contrast-enhanced imaging at 1--3 months is used to document complete necrosis.

\section*{Results and Interpretation}
A successful ablation is indicated by an avascular zone larger than the original tumor on follow-up imaging, with absence of enhancement. Persistent or rim-like enhancement at the ablation margin suggests residual viable tumor and may require repeat treatment or an alternative therapy.

\section*{Expected Results}
There is no normal result because ablation is a therapeutic intervention; the expected outcome is complete necrosis of the treated lesion. Success is defined by absence of residual enhancement at the site on follow-up imaging.

\section*{Follow-up}
Regular surveillance imaging (CT, MRI, or PET/CT) is performed at 1--3 months and then at intervals to detect local recurrence or new disease. Patients with metastatic disease usually continue systemic therapy under the direction of their oncology team.

\section*{References}
\begin{enumerate}
  \item MedlinePlus. Tumor Ablation. U.S. National Library of Medicine; 2025.
  \item Current Medical Diagnosis and Treatment. McGraw-Hill; 2025.
\end{enumerate}

\clearpage
